% RESEARCH_DRAFT_SHA256: 5f46bbea57ccb1c0996e432184beca3923262728d1865a1728323ad1d6c9507c
% Layout and source-reference wrappers expanded for the web renderer.
\documentclass[11pt]{article}
\usepackage[margin=25mm]{geometry}
\usepackage{amsthm}
\usepackage{fontspec,unicode-math}
\setmainfont{Latin Modern Roman}
\setmathfont{Latin Modern Math}
\usepackage{amsmath,amssymb,mathtools}
\usepackage{xurl,hyperref}
\hypersetup{colorlinks=true,urlcolor=blue,linkcolor=blue}
\setcounter{secnumdepth}{0}
\setlength{\parindent}{0pt}
\setlength{\parskip}{5pt}
\setlength{\emergencystretch}{3em}
\newcommand{\sref}[2]{\hyperlink{source:#1:#2}{[S#2]}}
\newcommand{\cataloguescope}[1]{\paragraph{Source question.}#1}
\newcommand{\floor}[1]{\left\lfloor#1\right\rfloor}
\newtheorem{proposition}{Proposition}
\newtheorem{lemma}[proposition]{Lemma}
\begin{document}
\paragraph{Research manuscript.} OpenAI \texttt{gpt-6-astra}, reasoning \texttt{ultra}; 11 October 2026.
\paragraph{Preserved source body.} The complete catalogue mathematical text and sources follow with layout and source links expanded, including source qualifications and its historical preparation notes. New research begins at the section headed Research attempt.

\section{Foundational Assumptions for Obfuscation, FHE, and SNARGs}\label{statement}

\subsection{Definitions and mathematical statement}

\paragraph{Shared notation.}Standard number systems are denoted by $\mathbb N,\mathbb Z,\mathbb Q,\mathbb R,\mathbb C$; index ranges, computational models and specialized notation follow the cited problem source and the definitions below.


Indistinguishability obfuscation hides which equivalent circuit was encoded. Virtual-black-box obfuscation compares information obtainable from the encoding with oracle access. Fully homomorphic encryption permits computation on encrypted inputs. SNARGs are succinct noninteractive arguments, with computational soundness.

A lattice generated by independent vectors $b_1,\ldots,b_d$ is $L=\{\sum_i z_i b_i:z_i\in\mathbb Z\}$. SVP seeks a shortest nonzero lattice vector; SIVP seeks $d$ independent short vectors, with approximation measured against the relevant successive minimum.

An $n$-qubit pure state is a unit vector in $(\mathbb C^2)^{\otimes n}$, considered up to global phase. A quantum algorithm applies permitted quantum operations and measurements; complexity depends on the gate and input-access model.

A cryptographic security requirement compares an efficient adversary\textquotesingle s advantage with a negligible function of the security parameter. A negligible function is smaller than every inverse polynomial eventually; computational and information-theoretic security are different requirements.
\paragraph{Statement recorded in the supplied source.}
Determine whether indistinguishability obfuscation can be based on weaker standard assumptions and made post-quantum secure; whether any meaningful virtual-black-box obfuscation can be constructed; whether fully homomorphic encryption can be based on objects other than integer lattices; and whether SNARGs for every NP language can be based on standard cryptographic assumptions. \href{https://doi.org/10.4171/owr/2025/3}{[S2]}
\paragraph{Scope qualification.}
The supplied assessment records impossibility of general virtual-black-box obfuscation and known obfuscation constructions, while leaving the stated weaker/post-quantum and other primitive questions unresolved. Meaningful restricted obfuscation is not equated with the impossible unrestricted version.
\paragraph{Source status.}
The supplied catalogue labels this record partially solved. Its source-status metadata was checked against the supplied snapshot on 2026-10-11. This import does not independently establish current open status or mathematical correctness. The catalogue literature assessment is dated 2026-08-22.
\subsection{Short English statement}
Resolve the bundled foundational and post-quantum construction questions for obfuscation, encrypted computation, and succinct arguments.
\subsection{Sources}
\begin{itemize}
\item[S1] Ulam AI and the UnsolvedMath Contributors, UnsolvedMath: A Curated Collection of Open Mathematics Problems, catalogue entry. Catalogue curation and metadata: CC BY 4.0. The supplied source question is retained; definitions and the English summary are editorial additions.
\url{https://huggingface.co/datasets/ulamai/UnsolvedMath}.
\item[S2] Cryptography (2025), Oberwolfach Reports, DOI 10.4171/owr/2025/3, pp. 141--174
\url{https://doi.org/10.4171/owr/2025/3}.
\item[S3] Barak et al., obfuscation impossibility
\url{https://people.seas.harvard.edu/~salil/research/obfuscate-abs.html}.
\item[S4] Jain et al., iO from well-founded assumptions
\url{https://doi.org/10.1145/3785007}.
\item[S5] OWR 2025/3 Cryptography report
\url{https://ems.press/content/serial-article-files/51349}.
\textit{Underlying Oberwolfach Reports material retains its source-specific licence, including CC BY-SA 4.0 where applicable; see THIRD\_PARTY\_NOTICES.md.}
\end{itemize}
\paragraph{Preparation and review.}
Original question attributed to the cited source and UnsolvedMath contributors. Codex (GPT-6; exact serving revision not exposed) assisted with definitions, English summaries, source inspection and formatting. No new human mathematical authorship or independent proof verification is asserted. Catalogue identity checked against the complete supplied registry and existing repository statements; known flagged repetitions and obvious identical conjectures excluded. Full mathematical-equivalence/literature review has not been performed. Statement preparation and final review followed contribution prompts C and D; automated checks verify structure and identity, not open status.
\label{end:statement}

\section{Research attempt}
\subsection{Literature, source context and scope}
The full original catalogue body and source qualifications are preserved above. The relevant primary workshop passages were read in the cached publisher report. Historical claims and subsequent attributed results are distinguished from the mathematical work below. A restricted example, standard lemma or known theorem is not counted as new progress on the entire original target. Missing assumptions and repeated or rhetorical extracts are not silently turned into different solved problems.

The source bundles distinct goals about obfuscation, post-quantum security, non-lattice fully homomorphic encryption and standard-assumption succinct arguments. It gives no single security experiment covering them. This attempt examines one tempting elementary substitute for the encryption component: perfectly secret additive masking with a short reusable key. It is a standard necessary-condition calculation, not a hardness-assumption reduction or an impossibility for computational FHE.
\subsection{A complete key-size obstruction under perfect secrecy}
Let a finite message space $\mathcal M$ and key space $\mathcal K$ support an encryption scheme, with key independent of the message, possibly randomized encryption, and ciphertexts that are finite binary strings (hence a countable discrete space). Suppose decryption from $(K,C)$ is always correct and the ciphertext distributions are identical for every plaintext.
\begin{proposition}
One has $|\mathcal K|\ge|\mathcal M|$. More generally, for any message distribution, perfect secrecy and correctness imply $H(K)\ge H(M)$, with entropy in bits.
\end{proposition}
\begin{proof}
Choose a ciphertext $c$ with positive probability for some plaintext. Perfect secrecy gives it positive probability for every plaintext $m$. For each such $m$ there is at least one key $k_m$ under which encryption of $m$ can yield $c$. Correct decryption of this same $c$ must return $m$ with that key. Two distinct plaintexts therefore cannot use the same selected key, proving the cardinality bound.

For the entropy statement, secrecy means $H(M\mid C)=H(M)$ and correctness means $H(M\mid K,C)=0$. Thus the chain-rule identity for conditional mutual information gives
\[
 H(M)=I(M;K\mid C)=H(K\mid C)-H(K\mid M,C)\le H(K\mid C)\le H(K).
\]
The inequalities use nonnegativity of conditional entropy and that conditioning does not increase entropy.
\end{proof}
For a uniformly random $t$-bit plaintext vector, this requires at least $t$ bits of key entropy. In particular it rules out perfect secrecy for arbitrarily many independent encrypted bits using a fixed-entropy reusable key.
\subsection{The additive-mask route fails on reuse}
A one-time pad $C=M\oplus K$ is perfectly secret for one bit when $K$ is uniform. Reusing that same key gives $C_1\oplus C_2=M_1\oplus M_2$. The two message vectors $(0,0)$ and $(0,1)$ produce ciphertext pairs of even and odd parity, respectively, so an observer distinguishes them without knowing the key. The ability to add ciphertexts does not itself provide multi-message privacy.
\section{Verification and remaining gaps}
Exhaustive one-bit tables verify single-use secrecy and disjoint parity supports on reuse. Computational security can evade the perfect-secrecy key bound, so this argument says nothing against existing or hypothetical non-lattice computational FHE. It also does not address the separate obfuscation or SNARG components, nor identify minimal assumptions or post-quantum reductions.

\section{References and source attribution}
Primary workshop report: \url{https://doi.org/10.4171/owr/2025/3}. Relevant passages read: OWR 2025/3, printed pages142--144, organizer questions on cryptographic primitives.

All retrievals were checked on 11 October 2026. Original question attribution remains in the preserved Sources section. External theorems are attributed and not claimed as new discoveries or fully verified merely by citation.


\paragraph{Integration review.} Codex (GPT-6; exact serving revision not exposed) reviewed the written arguments and expanded layout and source-reference wrappers. Repository metadata and the production web renderer were checked using Python 3.13 and JavaScript with MathJax 3.2.2. This model-assisted review is not independent human proof verification; no PDF compilation or formal proof checking is asserted.

\section{Completion Estimate}
\noindent\textbf{COMPLETION: 0\%}
No substantiated new progress toward the entire original target is claimed. The worked results and failed route are delimited above; a known result, illustrative model, or auxiliary obstruction is not counted as original-target completion.

\section{Final status and attribution}
\textbf{Attempt status: unresolved.} This describes this research attempt, not a new ruling on the source catalogue's status or on external literature claims.
The mathematics and this self-audit were generated by OpenAI \texttt{gpt-6-astra}, reasoning \texttt{ultra}, on 11 October 2026. Codex (GPT-6; exact serving revision not exposed) orchestrated the batch, read and reviewed the written mathematical arguments, independently reran the finite checks, and contributed formatting/repository review. No human mathematical authorship, independent human proof checking, or formal verification is claimed. The historical preparation paragraphs in the reproduced catalogue body concern that body only.

A separate OpenAI \texttt{gpt-6-astra}, reasoning \texttt{ultra}, invocation also read this written argument and independently performed the finite checks specified in the local peer-review audit. The exact serving revision is not exposed. This is model peer review, not independent human verification.

Contribution prompt A was applied to the complete definition and supplied source record, followed by the adversarial checks of prompt B. Local reproduction files are \texttt{batch026/checks.py}, \texttt{computations.json}, and the per-ID audit JSON. They were not public links at preparation time. Numerical checks are finite and do not prove unrestricted results. No LaTeX engine was found locally; compilation was not performed by this research agent.

\end{document}
